% ============================================================
%  Chapter 9 — Estimating Equations and Influence Functions
%  Compile standalone:  pdflatex chapter_9.tex
%  Compile via master:  pdflatex master.tex
% ============================================================
\documentclass[master.tex]{subfiles}

\begin{document}

% ------------------------------------------------------------
%  Title block (standalone) / chapter heading (master)
% ------------------------------------------------------------
\ifSubfilesClassLoaded{%
  \begin{center}
    {\LARGE\bfseries\color{isublue}
     Chapter 9: Estimating Equations and Influence Functions}\\[6pt]
    {\large Theory and Methods in Causal Inference}\\[4pt]
    {\normalsize Department of Statistics, Iowa State University}
  \end{center}
  \bigskip
  \hrule height 1.2pt \relax
  \smallskip
}{%
  \chapter{Estimating Equations and Influence Functions}
}

% ------------------------------------------------------------
%  Learning objectives
% ------------------------------------------------------------
\begin{tcolorbox}[defstyle, title={Learning Objectives}]
By the end of this lecture, students should be able to:
\begin{enumerate}
  \item Explain why identification of a causal parameter does not
    automatically yield a statistically reliable estimator, and
    articulate the distinct questions that an estimation theory must
    address.
  \item Define an estimating equation and verify the population moment
    condition for standard examples (sample mean, OLS, IPW).
  \item State the definition of an asymptotically linear estimator and
    derive the influence function from a generic first-order expansion
    of an estimating equation.
  \item Compute the influence function for simple causal estimators when
    nuisance functions are known, and interpret it as the
    infinitesimal contribution of a single observation.
  \item Describe how estimation error in nuisance functions propagates
    into the target estimator, and explain why this is the central
    statistical challenge in causal inference.
  \item Use the influence function to construct a consistent variance
    estimator and an approximate Wald confidence interval.
  \item Define efficiency in the class of regular estimators and explain
    the role of the efficient influence function as the foundation for
    doubly robust estimation in Chapter~10.
\end{enumerate}
\end{tcolorbox}


% ============================================================
\section{Why Estimation Needs Its Own Theory}
\label{w9:sec:motivation}
% ============================================================

So far the focus has been on \emph{identification}: under what
assumptions can a causal parameter be written as a functional of the
observed-data distribution?  Identification does not by itself provide
a statistically reliable estimator.  Once a causal parameter is
identified, several distinct questions remain:
\begin{itemize}
  \item How should the estimator be constructed from the observed sample?
  \item How does estimation of nuisance functions affect the target
    estimator?
  \item What is the large-sample distribution of the estimator?
  \item How can we compute valid standard errors and confidence intervals?
\end{itemize}

These questions motivate a separate theory of \emph{estimation and
inference}.  In causal inference this issue is especially important
because identified parameters often depend on auxiliary, or
\emph{nuisance}, functions such as the outcome regression
\[
  \mu_t(x) = \E(Y \mid T=t,\, X=x),
\]
or the propensity score
\[
  \pi(x) = P(T=1 \mid X=x).
\]
Even when the causal parameter is identified, different estimation
strategies may behave quite differently in terms of robustness,
efficiency, and sensitivity to model misspecification.

The goal of this chapter is to introduce a general framework for
estimation based on \emph{estimating equations} and \emph{influence
functions}.  This framework will serve as the foundation for the doubly
robust and semiparametric methods developed in Chapter~10; see also
\citet{imbens2015causal} and \citet{hernan2020causal} for complementary
treatments.

% ============================================================
\section{A Running Example: The Average Treatment Effect}
\label{w9:sec:ate}
% ============================================================

Throughout this chapter we use the average treatment effect (ATE) as a
running example.  Let
\[
  \psi = \E\{Y(1) - Y(0)\}.
\]
Under consistency, conditional exchangeability, and positivity, the ATE
is identified by the back-door formula (Chapter~5)
\[
  \psi = \E[\mu_1(X) - \mu_0(X)],
  \qquad
  \mu_t(x) = \E(Y \mid T=t,\, X=x),\quad t\in\{0,1\}.
\]

This identity tells us what the target parameter is, but it does not
uniquely determine how to estimate it.  For example, one may consider:
\begin{itemize}
  \item a \emph{regression-based} estimator obtained by fitting models
    for $\mu_1(x)$ and $\mu_0(x)$;
  \item a \emph{weighting} estimator based on the propensity score
    $\pi(x)$ (Chapter~6);
  \item an \emph{augmented} estimator combining both outcome and
    treatment models.
\end{itemize}
All of these estimators may target the same causal parameter, yet they
differ in their statistical properties.  A main goal of this chapter is
to develop a common language for describing and comparing such
estimators.

% ============================================================
\section{Estimating Equations}
\label{w9:sec:ee}
% ============================================================

A broad class of estimators can be defined as solutions to
\emph{estimating equations}.  Let $O_1,\dots,O_n$ be i.i.d.\
observations from a distribution $P$, and let $\theta$ denote a
finite-dimensional target parameter.  Write
\[
  \mathbb{P}_n f = \frac{1}{n}\sum_{i=1}^n f(O_i)
\]
for the empirical average.

\begin{defn}{Estimating Equation}
\label{w9:defn:ee}
An \textbf{estimating equation} for a parameter $\theta$ is an equation
of the form
\[
  \mathbb{P}_n\{U(O;\theta)\} = 0,
\]
where the \textbf{population moment condition}
\[
  \E\{U(O;\theta_0)\} = 0
\]
holds at the true parameter value $\theta_0$.  The function
$U(O;\theta)$ is called the \textbf{estimating function}.
\end{defn}

An estimating-equation estimator $\hat\theta$ is any solution to the
sample moment condition above.  Many familiar estimators take this form.

\begin{example}{Sample mean}
\label{w9:ex:mean}
Let $\theta = \E(Y)$.  The sample mean $\hat\theta = \bar Y$ solves
\[
  \mathbb{P}_n(Y-\theta) = 0.
\]
The estimating function is $U(O;\theta) = Y - \theta$.
\end{example}

\begin{example}{Ordinary least squares}
\label{w9:ex:ols}
Let $O = (X,Y)$ and suppose $\beta$ is defined by the linear moment
condition $\E[X\{Y - X^\top\beta\}] = 0$.  Then the OLS estimator
solves
\[
  \mathbb{P}_n\bigl[X\{Y - X^\top\beta\}\bigr] = 0.
\]
\end{example}

\begin{example}{A simple IPW estimating equation}
\label{w9:ex:ipw-ee}
Suppose the propensity score $\pi(X)$ is known and consider the ATE
parameter
\[
  \psi = \E\left\{\frac{TY}{\pi(X)} - \frac{(1-T)Y}{1-\pi(X)}\right\}.
\]
Then $\psi$ solves
\[
  \E\!\left[\frac{TY}{\pi(X)} - \frac{(1-T)Y}{1-\pi(X)} - \psi\right] = 0,
\]
and the corresponding estimator replaces the expectation by the
empirical average.
\end{example}

The main advantage of the estimating-equation framework is that it
provides a unified language for both estimator construction and
asymptotic analysis.

% ============================================================
\section{From Estimating Equations to Asymptotic Linearity}
\label{w9:sec:asymlin}
% ============================================================

Estimating equations are useful not only because they define estimators,
but also because they often yield a convenient first-order expansion.
Under suitable regularity conditions, an estimator solving an
estimating equation can typically be approximated as
\[
  \sqrt{n}(\hat\theta - \theta_0)
  =
  \frac{1}{\sqrt{n}}\sum_{i=1}^n \varphi(O_i) + o_p(1)
\]
for some mean-zero function $\varphi(O)$.

\begin{defn}{Asymptotic Linearity and Influence Function}
\label{w9:defn:asymlin}
An estimator $\hat\theta$ is \textbf{asymptotically linear} with
\textbf{influence function} $\varphi(O)$ if
\[
  \sqrt{n}(\hat\theta - \theta_0)
  =
  \frac{1}{\sqrt{n}}\sum_{i=1}^n \varphi(O_i) + o_p(1),
  \qquad
  \E\{\varphi(O)\} = 0.
\]
\end{defn}

This representation is fundamental because it immediately implies
asymptotic normality by the multivariate central limit theorem.  When
$\theta \in \mathbb{R}^p$ the influence function $\varphi(O)$ is also
$\mathbb{R}^p$-valued, and
\[
  \sqrt{n}(\hat\theta - \theta_0)
  \overset{d}{\longrightarrow}
  N\!\left(\mathbf{0},\; \Sigma\right),
  \qquad
  \Sigma = \E\bigl[\varphi(O)\varphi(O)^\top\bigr],
\]
where $\Sigma$ is the $p \times p$ asymptotic covariance matrix.  In
the scalar case $p = 1$ this reduces to $\Sigma = \E[\varphi(O)^2]$.
Once the influence function $\varphi(O)$ is identified, we obtain both
the asymptotic distribution and a principled basis for variance
estimation and confidence intervals.

\begin{prop}{Generic first-order expansion}
\label{w9:prop:expansion}
Under smoothness and regularity conditions, an estimator $\hat\theta$
solving $\mathbb{P}_n\{U(O;\theta)\}=0$ admits a first-order
expansion
\[
  \sqrt{n}(\hat\theta - \theta_0)
  =
  -A^{-1}\,\frac{1}{\sqrt{n}}\sum_{i=1}^n U(O_i;\theta_0)
  + o_p(1),
\]
where
\[
  A = \E\!\left\{\frac{\partial}{\partial\theta^\top}U(O;\theta_0)\right\}.
\]
Hence the influence function is $\varphi(O) = -A^{-1}U(O;\theta_0)$.
\end{prop}

\begin{proof}
A Taylor expansion of the sample moment condition around $\theta_0$
gives
\[
  0 = \mathbb{P}_n\{U(O;\hat\theta)\}
    \approx
    \mathbb{P}_n\{U(O;\theta_0)\}
    + \mathbb{P}_n\!\left\{\frac{\partial}{\partial\theta^\top}
      U(O;\theta_0)\right\}(\hat\theta - \theta_0).
\]
By the law of large numbers,
$\mathbb{P}_n\{\partial_{\theta^\top} U(O;\theta_0)\} \to A$ in
probability.  Rearranging and multiplying by $\sqrt{n}$ yields the
stated expansion.
\end{proof}

Proposition~\ref{w9:prop:expansion} is a standard M-estimation result
and shows why estimating equations naturally lead to asymptotic
linearity.

% ============================================================
\section{Influence Functions: Intuition}
\label{w9:sec:if-intuition}
% ============================================================

\subsection{Statistical Functionals}
\label{w9:subsec:functional}

Many parameters of interest in statistics and causal inference can be
written as \emph{functionals} of the underlying distribution.  We make
this precise before defining the influence function.

\begin{defn}{Statistical Functional}
\label{w9:defn:functional}
Let $\mathcal{P}$ be a collection of probability distributions on the
sample space of $O$.  A \textbf{statistical functional} is a map
\[
  \Psi : \mathcal{P} \longrightarrow \mathbb{R}^p
\]
that assigns a parameter value $\psi = \Psi(P)$ to each distribution
$P \in \mathcal{P}$.  The \textbf{plug-in estimator} of $\psi$ is
$\hat\psi = \Psi(\mathbb{P}_n)$, obtained by replacing $P$ with the
empirical distribution $\mathbb{P}_n$.
\end{defn}

Most parameters encountered in this course are statistical functionals.

\begin{example}{Common statistical functionals}
\label{w9:ex:functionals}
\begin{enumerate}[label=(\roman*)]
  \item \textbf{Mean.}
    $\Psi(P) = \E_P(Y) = \int y\,dP(y)$.
    The plug-in estimator is $\bar Y$.
  \item \textbf{Variance.}
    $\Psi(P) = \E_P(Y^2) - [\E_P(Y)]^2$.
    The plug-in estimator is the sample variance (with $n$ in the
    denominator).
  \item \textbf{Quantile.}
    $\Psi(P) = F_P^{-1}(\tau) = \inf\{y : F_P(y) \geq \tau\}$ for
    $\tau \in (0,1)$.  The plug-in estimator is the sample
    $\tau$-quantile.
  \item \textbf{OLS regression coefficient.}
    $\Psi(P) = [\E_P(XX^\top)]^{-1}\E_P(XY)$.
    The plug-in estimator is the OLS estimator
    $(\sum_i X_iX_i^\top)^{-1}\sum_i X_iY_i$.
  \item \textbf{Average treatment effect.}
    $\Psi(P) = \E_P[\mu_1(X) - \mu_0(X)]$ under the back-door
    identification formula.  The plug-in estimator substitutes
    estimated regression functions $\hat\mu_t$.
\end{enumerate}
\end{example}

\begin{remark}[Linearity vs.\ nonlinearity]
\label{w9:rem:linear-functional}
A functional $\Psi$ is \textbf{linear} if
$\Psi((1-\epsilon)P + \epsilon Q) = (1-\epsilon)\Psi(P) +
\epsilon\Psi(Q)$ for all $P, Q \in \mathcal{P}$ and
$\epsilon \in [0,1]$.  The mean (i) is linear; the variance (ii),
quantile (iii), OLS coefficient (iv), and ATE (v) are nonlinear.
Linear functionals are straightforward to analyze because their
plug-in estimators are sample averages.  Nonlinear functionals require
a local linearization, which is precisely what the influence function
provides.
\end{remark}

\subsection{Influence Functions}
\label{w9:subsec:if-def}

Influence functions describe the first-order sensitivity of a
statistical functional to small perturbations of the underlying
distribution.  Let $\Psi(P)$ be a parameter of interest viewed as a
functional of the observed-data distribution $P$.  Informally, the
influence function measures how $\Psi(P)$ changes when $P$ is slightly
contaminated in the direction of a point mass at a single observation
$O$.

\begin{quote}
\itshape
The influence function is the infinitesimal contribution of a single
observation to the first-order behavior of the estimator.
\end{quote}

\begin{defn}{Influence Function (informal)}
\label{w9:defn:if}
Let $\psi = \Psi(P)$ be a scalar parameter.  A function $\varphi(O)$
is called an \textbf{influence function} for $\Psi$ if it has mean zero
under $P$ and gives the first-order derivative of $\Psi(P)$ along
regular parametric submodels, i.e.,
\[
  \hat\psi - \psi
  =
  \frac{1}{n}\sum_{i=1}^n \varphi(O_i) + o_p(n^{-1/2}).
\]
\end{defn}

\begin{remark}[Formal derivation via pathwise derivatives]
\label{w9:rem:pathwise}
The influence function can be defined formally through the
\emph{pathwise (Gateaux) derivative} of the functional $\Psi(P)$: if
$P_\epsilon = (1-\epsilon)P + \epsilon\delta_O$ is the $\epsilon$-mixture
of $P$ with a point mass at $O$, then the influence function satisfies
\[
  \varphi(O) = \left.\frac{d}{d\epsilon}\Psi(P_\epsilon)\right|_{\epsilon=0}.
\]
This derivative-based definition is the starting point for
semiparametric efficiency theory.  At the level of this course it is
enough to interpret $\varphi(O)$ as the object appearing in the
asymptotic linear expansion, as in Definition~\ref{w9:defn:asymlin}.
\end{remark}

\begin{example}{Influence function of the OLS slope via pathwise derivative}
\label{w9:ex:if-ols}
Let $O = (X, Y) \in \mathbb{R}^{p+1}$ and consider the population
linear regression coefficient
\[
  \beta = \Psi(P)
  =
  \bigl[\E(XX^\top)\bigr]^{-1}\E(XY),
\]
a nonlinear functional of $P$.  The OLS estimator
$\hat\beta = \bigl(\sum_i X_iX_i^\top\bigr)^{-1}\sum_i X_iY_i$ is
the empirical plug-in $\Psi(\mathbb{P}_n)$.

\medskip\noindent\textbf{Step 1: Construct the perturbed distribution.}
Following Remark~\ref{w9:rem:pathwise}, let
$P_\epsilon = (1-\epsilon)P + \epsilon\delta_O$ be the mixture of $P$
with a point mass at a fixed observation $O = (X, Y)$.  Then
\begin{align*}
  \E_{P_\epsilon}(\tilde X \tilde X^\top)
  &= (1-\epsilon)\E(XX^\top) + \epsilon XX^\top, \\
  \E_{P_\epsilon}(\tilde X \tilde Y)
  &= (1-\epsilon)\E(XY) + \epsilon XY,
\end{align*}
so the functional evaluated at $P_\epsilon$ is
\[
  \Psi(P_\epsilon)
  =
  \bigl[(1-\epsilon)\E(XX^\top) + \epsilon XX^\top\bigr]^{-1}
  \bigl[(1-\epsilon)\E(XY) + \epsilon XY\bigr].
\]

\medskip\noindent\textbf{Step 2: Differentiate with respect to
$\epsilon$ at $\epsilon = 0$.}
Write $M(\epsilon) = (1-\epsilon)\E(XX^\top) + \epsilon XX^\top$ and
$v(\epsilon) = (1-\epsilon)\E(XY) + \epsilon XY$, so
$\Psi(P_\epsilon) = M(\epsilon)^{-1}v(\epsilon)$.  Note that both
paths are linear in $\epsilon$, with derivatives
\[
  \dot M(0)
  \;:=\;
  \left.\frac{d}{d\epsilon}M(\epsilon)\right|_{\epsilon=0}
  = XX^\top - \E(XX^\top),
  \qquad
  \dot v(0)
  \;:=\;
  \left.\frac{d}{d\epsilon}v(\epsilon)\right|_{\epsilon=0}
  = XY - \E(XY).
\]

\medskip\noindent\textit{Deriving the derivative of a matrix inverse.}
For any differentiable matrix-valued path $M(\epsilon)$ with $M(\epsilon)$
invertible, differentiating the identity
$M(\epsilon)M(\epsilon)^{-1} = I$ with respect to $\epsilon$ gives
\[
  \dot M(\epsilon)\,M(\epsilon)^{-1}
  + M(\epsilon)\,\frac{d}{d\epsilon}M(\epsilon)^{-1}
  = 0.
\]
Solving for the derivative of the inverse,
\[
  \frac{d}{d\epsilon}M(\epsilon)^{-1}
  =
  -M(\epsilon)^{-1}\dot M(\epsilon)\,M(\epsilon)^{-1}.
\]
Evaluating at $\epsilon = 0$, with $M(0) = \E(XX^\top)$,
\[
  \left.\frac{d}{d\epsilon}M(\epsilon)^{-1}\right|_{\epsilon=0}
  =
  -\bigl[\E(XX^\top)\bigr]^{-1}
  \bigl(XX^\top - \E(XX^\top)\bigr)
  \bigl[\E(XX^\top)\bigr]^{-1}.
\]

\medskip\noindent\textit{Applying the product rule.}
Using $\frac{d}{d\epsilon}[M^{-1}v] = \dot{(M^{-1})}v + M^{-1}\dot v$,
\begin{align*}
  \varphi(O)
  &= \left.\frac{d}{d\epsilon}\Psi(P_\epsilon)\right|_{\epsilon=0} \\
  &= -\bigl[\E(XX^\top)\bigr]^{-1}
     \bigl(XX^\top - \E(XX^\top)\bigr)
     \bigl[\E(XX^\top)\bigr]^{-1}\E(XY)
     +
     \bigl[\E(XX^\top)\bigr]^{-1}\bigl(XY - \E(XY)\bigr).
\end{align*}
Substituting $\beta = [\E(XX^\top)]^{-1}\E(XY)$ and collecting terms,
\begin{align*}
  \varphi(O)
  &= \bigl[\E(XX^\top)\bigr]^{-1}
     \bigl\{XY - \E(XY) - (XX^\top - \E(XX^\top))\beta\bigr\} \\
  &= \bigl[\E(XX^\top)\bigr]^{-1}
     X\bigl(Y - X^\top\beta\bigr)
     - \bigl[\E(XX^\top)\bigr]^{-1}
     \bigl\{\E(XY) - \E(XX^\top)\beta\bigr\}.
\end{align*}
The second term is zero because $\beta = [\E(XX^\top)]^{-1}\E(XY)$
implies $\E(XY) - \E(XX^\top)\beta = 0$.  Hence
\[
  \boxed{\varphi(O)
  =
  \bigl[\E(XX^\top)\bigr]^{-1}X\bigl(Y - X^\top\beta\bigr).}
\]

\medskip\noindent\textbf{Verification.}
One can confirm $\E\{\varphi(O)\} = 0$ directly:
$\E[X(Y - X^\top\beta)] = \E(XY) - \E(XX^\top)\beta = 0$ by
definition of $\beta$.

\medskip\noindent\textbf{Connection to the estimating-equation
approach.}  The OLS estimating function is
$U(O;\beta) = X(Y - X^\top\beta)$, so
$A = \E\{\partial_{\beta^\top}U\} = -\E(XX^\top)$ and
Proposition~\ref{w9:prop:expansion} gives
$\varphi(O) = -A^{-1}U(O;\beta_0) = [\E(XX^\top)]^{-1}X(Y -
X^\top\beta)$, agreeing with the pathwise calculation.  The two
approaches yield the same answer; the pathwise route is more
fundamental, while the estimating-equation route is more
computationally direct.
\end{example}

% ============================================================
\section{Influence Functions for Simple Causal Estimators}
\label{w9:sec:if-causal}
% ============================================================

We now illustrate the idea in causal settings.

\subsection{Known outcome regression}
\label{w9:subsec:if-reg}

Suppose
\[
  \psi = \E[\mu_1(X) - \mu_0(X)]
\]
and the functions $\mu_1(\cdot)$ and $\mu_0(\cdot)$ are known.  The
natural plug-in estimator is
\[
  \hat\psi = \mathbb{P}_n\{\mu_1(X) - \mu_0(X)\}.
\]
Its influence function is
\[
  \varphi(O) = \mu_1(X) - \mu_0(X) - \psi.
\]

\subsection{Known propensity score}
\label{w9:subsec:if-ipw}

Suppose $\pi(X)$ is known and
\[
  \psi = \E\!\left\{\frac{TY}{\pi(X)} - \frac{(1-T)Y}{1-\pi(X)}\right\}.
\]
The empirical analogue is
\[
  \hat\psi
  =
  \mathbb{P}_n\!\left\{\frac{TY}{\pi(X)} - \frac{(1-T)Y}{1-\pi(X)}\right\},
\]
and its influence function is
\[
  \varphi(O)
  =
  \frac{TY}{\pi(X)} - \frac{(1-T)Y}{1-\pi(X)} - \psi.
\]

\begin{remark}[Nuisance functions and the influence function]
\label{w9:rem:if-principle}
These two examples illustrate an important principle: when nuisance
functions are \emph{known}, the influence function is simply the
centered estimating function.  This observation becomes much more
consequential in the next chapter, where nuisance functions are
unknown and must be estimated from data.  The key challenge will be to
understand how estimation error in $\hat\mu_t$ or $\hat\pi$ perturbs
the first-order expansion.
\end{remark}

\begin{remark}[Estimating function, influence function, and efficient influence function]
\label{w9:rem:three-objects}
These three objects are closely related, but they are not the same.

\begin{enumerate}
  \item An \emph{estimating function} defines an estimator through a
    sample moment equation such as
    \[
      \mathbb{P}_n\{\phi(O;\theta)\} = 0.
    \]

  \item An \emph{influence function} describes the first-order
    asymptotic behavior of an estimator:
    \[
      \sqrt{n}(\hat\theta - \theta_0)
      =
      \frac{1}{\sqrt{n}}\sum_{i=1}^n \varphi(O_i) + o_p(1).
    \]

  \item The \emph{efficient influence function} is the influence
    function with the smallest possible asymptotic variance among
    regular estimators in the statistical model under consideration.
\end{enumerate}

In simple examples with known nuisance functions, the centered
estimating function often coincides with the influence function.
However, this is not automatic in general, especially when nuisance
parameters are estimated (Section~\ref{w9:sec:nuisance}).  The
efficient influence function is a further refinement: it is not merely
a valid influence function, but the one corresponding to the
semiparametric efficiency bound (Section~\ref{w9:subsec:eif}).
\end{remark}

% ============================================================
\section{Z-Estimation with Nuisance Parameters}
\label{w9:sec:nuisance}
% ============================================================

The examples in Section~\ref{w9:sec:if-causal} treated nuisance
functions such as $\mu_t(x)$ and $\pi(x)$ as known.  That
simplification was pedagogically useful because it made the influence
function transparent, but it is not realistic in practice.  In real
applications, nuisance functions must be estimated from the same data
used to estimate the target parameter.  Once this happens, the
first-order expansion of $\hat\psi$ is no longer obtained by simply
centering the estimating function with the nuisance held fixed: the
randomness in the estimated nuisance introduces an additional term that
the naive plug-in formula ignores.  We therefore need a framework that
accounts jointly for estimation of the target parameter and the nuisance
parameter.  Z-estimation provides that framework: it treats both as
components of a single stacked estimating equation, so the corrected
influence function --- including the nuisance-estimation term --- falls
out of one unified first-order expansion.

% ------------------------------------------------------------
\subsection{The Stacked Estimating Equation Framework}
\label{w9:subsec:stacked}
% ------------------------------------------------------------

The key idea is that once $\alpha$ is estimated rather than known, the
target equation for $\psi$ and the nuisance equation for $\alpha$ must
be analyzed jointly rather than one after the other: only the joint
analysis reveals the additional term that the estimated nuisance
contributes to the first-order expansion of $\hat\psi$.

Suppose the target parameter $\psi \in \mathbb{R}$ depends on an
unknown nuisance parameter $\alpha \in \mathbb{R}^q$ (e.g., regression
coefficients in a parametric model for $\pi(x)$).  Assume both are
estimated by solving a \emph{stacked} system of estimating equations:
\begin{align}
  \mathbb{P}_n\bigl\{U_1(O;\psi,\alpha)\bigr\} &= 0,
  \label{w9:eq:ee-target}\\
  \mathbb{P}_n\bigl\{U_2(O;\alpha)\bigr\} &= 0.
  \label{w9:eq:ee-nuisance}
\end{align}
Here $U_1(O;\psi,\alpha)$ is scalar and defines the target estimator
$\hat\psi$ given $\alpha$, while $U_2(O;\alpha) \in \mathbb{R}^q$
defines the nuisance estimator $\hat\alpha$.  The estimating equations
\eqref{w9:eq:ee-target}--\eqref{w9:eq:ee-nuisance} are solved
\emph{simultaneously} (or sequentially: first \eqref{w9:eq:ee-nuisance}
for $\hat\alpha$, then \eqref{w9:eq:ee-target} for $\hat\psi$).

\begin{remark}[Why stack the equations?]
\label{w9:rem:why-stack}
One might think it is enough to plug the estimated $\hat\alpha$ into
the influence function derived under known $\alpha$ and proceed as in
Section~\ref{w9:sec:if-causal}.  This is incorrect in general: the
randomness in $\hat\alpha$ contributes an additional term to the
first-order expansion of $\hat\psi$ that is not captured by the
plug-in influence function.  The stacked framework keeps track of this
extra term automatically.
\end{remark}

\subsection{The Z-Estimation Theorem}
\label{w9:subsec:zthm}

Write $\theta = (\psi^\top, \alpha^\top)^\top \in \mathbb{R}^{p+q}$,
where $\psi \in \mathbb{R}^p$ is the target parameter and
$\alpha \in \mathbb{R}^q$ is the nuisance parameter, and let
\[
  U(O;\theta)
  =
  \begin{pmatrix} U_1(O;\psi,\alpha) \\ U_2(O;\alpha) \end{pmatrix}
\]
be the stacked estimating function.  The joint estimator
$\hat\theta = (\hat\psi, \hat\alpha^\top)^\top$ solves
$\mathbb{P}_n\{U(O;\hat\theta)\} = 0$.

\begin{thm}{Z-Estimation Theorem (Stacked Equations)}
\label{w9:thm:zest}
Suppose the true parameter value is $\theta_0 = (\psi_0,\alpha_0^\top)^\top$
and the following regularity conditions hold:
\begin{enumerate}[label=(\roman*)]
  \item $\E\{U(O;\theta_0)\}=0$ (population moment condition);
  \item the Jacobian
    $\mathcal{A} = \E\{\partial_{\theta^\top} U(O;\theta_0)\}$ is
    invertible;
  \item $U(O;\theta)$ is smooth enough in $\theta$ that a
    uniform law of large numbers and a CLT apply.
\end{enumerate}
Then $\hat\theta$ is consistent for $\theta_0$ and
\[
  \sqrt{n}(\hat\theta - \theta_0)
  =
  -\mathcal{A}^{-1}\,
  \frac{1}{\sqrt{n}}\sum_{i=1}^n U(O_i;\theta_0)
  + o_p(1).
\]
In particular, after partitioning $\mathcal{A}^{-1}$ conformably with
$(\psi,\alpha)$ as detailed below, the influence function of $\hat\psi$
is
\[
  \varphi(O)
  =
  -A_{11}^{-1}\,U_1(O;\psi_0,\alpha_0)
  + A_{11}^{-1}A_{12}A_{22}^{-1}\,U_2(O;\alpha_0),
\]
where $A_{11}$, $A_{12}$, $A_{22}$ are the blocks of $\mathcal{A}$
defined in \eqref{w9:eq:if-corrected} below.
\end{thm}

\begin{proof}
The argument mirrors Proposition~\ref{w9:prop:expansion} applied to the
full stacked system.  A Taylor expansion of
$\mathbb{P}_n\{U(O;\hat\theta)\}=0$ around $\theta_0$ gives
\[
  0
  \approx
  \mathbb{P}_n\{U(O;\theta_0)\}
  +
  \mathbb{P}_n\!\left\{\frac{\partial}{\partial\theta^\top}
  U(O;\theta_0)\right\}(\hat\theta-\theta_0).
\]
By the law of large numbers, the matrix of derivatives converges to
$\mathcal{A}$.  Rearranging and multiplying by $\sqrt{n}$,
\[
  \sqrt{n}(\hat\theta - \theta_0)
  \approx
  -\mathcal{A}^{-1}\,\frac{1}{\sqrt{n}}\sum_{i=1}^n U(O_i;\theta_0).
\]
Extracting the first component gives the influence function of
$\hat\psi$.
\end{proof}

To see the structure explicitly, partition $\mathcal{A}$ conformably
with $(\psi,\alpha)$:
\[
  \mathcal{A}
  =
  \begin{pmatrix} A_{11} & A_{12} \\ 0 & A_{22} \end{pmatrix},
\]
where $A_{11} = \E\{\partial_\psi U_1\}$,
$A_{12} = \E\{\partial_{\alpha^\top} U_1\}$, and
$A_{22} = \E\{\partial_{\alpha^\top} U_2\}$.
The lower-left block is zero because $U_2$ does not depend on $\psi$.
By block-matrix inversion,
\[
  \mathcal{A}^{-1}
  =
  \begin{pmatrix}
    A_{11}^{-1} & -A_{11}^{-1}A_{12}A_{22}^{-1} \\
    0 & A_{22}^{-1}
  \end{pmatrix}.
\]
Reading off the first row, the influence function of $\hat\psi$ is
\begin{equation}
  \varphi(O)
  =
  -A_{11}^{-1}\,U_1(O;\psi_0,\alpha_0)
  +A_{11}^{-1}A_{12}A_{22}^{-1}\,U_2(O;\alpha_0).
  \label{w9:eq:if-corrected}
\end{equation}

\begin{remark}[Interpreting the correction term]
\label{w9:rem:correction}
Equation~\eqref{w9:eq:if-corrected} decomposes the influence function
into two parts.  The first term,
$-A_{11}^{-1}U_1(O;\psi_0,\alpha_0)$, is exactly what we would obtain
if $\alpha_0$ were known.  The second term,
$A_{11}^{-1}A_{12}A_{22}^{-1}U_2(O;\alpha_0)$, is a
\emph{nuisance correction} that accounts for the additional variability
introduced by estimating $\alpha$.  If $A_{12}=0$, the target estimating
function $U_1$ does not depend on $\alpha$ at $\alpha_0$, and the
correction vanishes: knowing $\hat\alpha$ confers no first-order
benefit or cost.  This is the key condition exploited by
\emph{locally efficient} and \emph{doubly robust} estimators in
Chapter~10.
\end{remark}

% ------------------------------------------------------------
\subsection{Working Example: IPW with Estimated Propensity Score}
\label{w9:subsec:ipw-example}
% ------------------------------------------------------------

Chapter~6 introduced IPW as an estimator motivated by the back-door
identification formula and the propensity score.  The present example
revisits the same estimator from a different angle: not identification,
but first-order asymptotic analysis when the propensity score is itself
estimated.  Concretely, we apply the Z-estimation theorem to the IPW
estimator of the ATE when the propensity score is estimated by logistic
regression.

\medskip\noindent\textbf{Setup.}
Assume a logistic regression model $\pi(X;\alpha) =
\mathrm{expit}(X^\top\alpha)$, where
$\mathrm{expit}(u) = e^u/(1+e^u)$.
The IPW estimating equation for $\psi$ is
\[
  U_1(O;\psi,\alpha)
  =
  \frac{TY}{\pi(X;\alpha)}
  - \frac{(1-T)Y}{1-\pi(X;\alpha)}
  - \psi.
\]
The score equation for the logistic regression is
\[
  U_2(O;\alpha)
  =
  X\bigl\{T - \pi(X;\alpha)\bigr\}.
\]
The stacked system $\mathbb{P}_n U_1 = 0$, $\mathbb{P}_n U_2 = 0$
yields $(\hat\psi, \hat\alpha)$.

\medskip\noindent\textbf{Computing the Jacobian blocks.}
\begin{align*}
  A_{11}
  &= \E\!\left\{\frac{\partial}{\partial\psi}U_1\right\} = -1.\\[4pt]
  A_{22}
  &= \E\!\left\{\frac{\partial}{\partial\alpha^\top}U_2\right\}
   = -\E\!\left\{\pi(X;\alpha_0)\bigl(1-\pi(X;\alpha_0)\bigr)
                 XX^\top\right\}
   =: -\Sigma_\pi.\\[4pt]
  A_{12}
  &= \E\!\left\{\frac{\partial}{\partial\alpha^\top}U_1\right\}
   = \E\!\left\{Y\left(
       \frac{T}{\pi^2}\,\pi(1-\pi)X^\top
      +\frac{1-T}{(1-\pi)^2}\,\pi(1-\pi)X^\top
     \right)\right\}\\
  &= \E\!\left\{\frac{Y\,X^\top}{\pi(X;\alpha_0)
       \bigl(1-\pi(X;\alpha_0)\bigr)}
     \bigl\{T(1-\pi) + (1-T)\pi\bigr\}\right\}\\
  &= \E\!\left\{Y X^\top
     \left(\frac{T}{\pi(X;\alpha_0)}
           - \frac{1-T}{1-\pi(X;\alpha_0)}\right)
     \cdot\frac{\pi(X;\alpha_0)(1-\pi(X;\alpha_0))}
               {\pi(X;\alpha_0)(1-\pi(X;\alpha_0))}\right\}.
\end{align*}
A cleaner expression is obtained by noting that
\[
  \frac{\partial}{\partial\alpha^\top}
  \left\{\frac{T}{\pi(X;\alpha)}-\frac{1-T}{1-\pi(X;\alpha)}\right\}
  =
  -\frac{T}{\pi^2}\,\dot\pi X^\top
  +\frac{1-T}{(1-\pi)^2}\,\dot\pi X^\top,
\]
where $\dot\pi = \pi(1-\pi)$.  Evaluating at $\alpha_0$ and taking
expectations,
\[
  A_{12}
  =
  \E\!\left[Y\left\{
    -\frac{T}{\pi(X;\alpha_0)}
    +\frac{1-T}{1-\pi(X;\alpha_0)}
  \right\}X^\top\right].
\]

\medskip\noindent\textbf{The corrected influence function.}
Substituting into \eqref{w9:eq:if-corrected} with $A_{11}=-1$ and
$A_{22}^{-1} = -\Sigma_\pi^{-1}$,
\begin{align}
  \varphi(O)
  &=
  U_1(O;\psi_0,\alpha_0)
  - A_{12}\,\Sigma_\pi^{-1}\,U_2(O;\alpha_0).
  \label{w9:eq:if-ipw-corrected}
\end{align}
The first term is the naive IPW influence function from
Section~\ref{w9:subsec:if-ipw}.  The second term is the nuisance
correction: it removes the first-order contribution of
$\hat\alpha - \alpha_0$ to the estimation error in $\hat\psi$.

\begin{example}{Variance of the IPW estimator with estimated propensity score}
\label{w9:ex:ipw-var}
We show that estimating $\alpha$ by maximum likelihood (ML) can only
\emph{reduce} the asymptotic variance of $\hat\psi_{\mathrm{IPW}}$
relative to using the true $\alpha_0$.  Two Bartlett-type identities
are the key.

\medskip\noindent\textbf{Asymptotic variance with known vs.\ estimated
propensity score.}
Write $\sigma_1^2 = \E[U_1(O;\psi_0,\alpha_0)^2]$ for the variance of
the naive IPW estimating function when $\alpha_0$ is treated as known.
From the corrected influence function \eqref{w9:eq:if-ipw-corrected}
with $A_{11} = -1$, the asymptotic variance with estimated $\hat\alpha$
is
\begin{align}
  \E[\varphi(O)^2]
  &= \E\!\left[
      \bigl(U_1 - A_{12}\Sigma_\pi^{-1}U_2\bigr)^2
    \right] \notag\\
  &= \sigma_1^2
    - 2\,A_{12}\Sigma_\pi^{-1}\E[U_1 U_2^\top]
    + A_{12}\Sigma_\pi^{-1}\E[U_2 U_2^\top]\Sigma_\pi^{-1}A_{12}^\top.
  \label{w9:eq:var-decomp}
\end{align}

\medskip\noindent\textbf{Identity 1: Bartlett identity for the ML
score.}
Because $U_2(O;\alpha) = X\{T - \pi(X;\alpha)\}$ is the ML score for
$\alpha$ in a correctly specified logistic model, it satisfies the
Bartlett identity (information equality):
\[
  \E\!\left[\frac{\partial}{\partial\alpha^\top}U_2(O;\alpha_0)\right]
  +
  \E\bigl[U_2(O;\alpha_0)\,U_2(O;\alpha_0)^\top\bigr]
  = 0.
\]
Since $A_{22} = \E\{\partial_{\alpha^\top}U_2\} = -\Sigma_\pi$, this
gives
\begin{equation}
  \E\bigl[U_2 U_2^\top\bigr] = \Sigma_\pi.
  \label{w9:eq:bartlett}
\end{equation}

\medskip\noindent\textbf{Identity 2: the analogous identity for $U_1$.}
A parallel identity holds for $U_1$.  Because the marginal
distribution of $X$ does not depend on $\alpha$, the IPW estimating
function satisfies
\[
  \E_\alpha[U_1(O;\psi_0,\alpha)]
  =
  \E_X[\mu_1(X)] - \E_X[\mu_0(X)] - \psi_0
  = 0
  \quad\text{for all }\alpha,
\]
where the expectation $\E_\alpha$ is under the model with propensity
$\pi(X;\alpha)$.  Differentiating this identity with respect to
$\alpha^\top$ and applying the Leibniz rule,
\[
  0
  =
  \frac{\partial}{\partial\alpha^\top}\E_\alpha[U_1]
  =
  \E_{\alpha}\!\left[\frac{\partial U_1}{\partial\alpha^\top}\right]
  +
  \E_{\alpha}[U_1 U_2^\top].
\]
Evaluating at $\alpha = \alpha_0$ gives
\begin{equation}
  A_{12} = -\E[U_1 U_2^\top].
  \label{w9:eq:a12-identity}
\end{equation}

\medskip\noindent\textbf{Variance reduction.}
Substituting \eqref{w9:eq:bartlett} and \eqref{w9:eq:a12-identity}
into \eqref{w9:eq:var-decomp},
\begin{align*}
  \E[\varphi(O)^2]
  &= \sigma_1^2
     + 2\,\E[U_1 U_2^\top]\Sigma_\pi^{-1}\E[U_1 U_2^\top]
     - \E[U_1 U_2^\top]\Sigma_\pi^{-1}\Sigma_\pi
       \Sigma_\pi^{-1}\E[U_2 U_1^\top] \\
  &= \sigma_1^2
     - \E[U_1 U_2^\top]\Sigma_\pi^{-1}\E[U_2 U_1^\top].
\end{align*}
Since $\Sigma_\pi$ is positive definite,
$\E[U_1 U_2^\top]\Sigma_\pi^{-1}\E[U_2 U_1^\top] \geq 0$, and
therefore
\[
  \E[\varphi(O)^2] \leq \sigma_1^2,
\]
with equality only when $\E[U_1 U_2^\top] = 0$.  Thus \emph{ML
estimation of the propensity score parameters reduces the asymptotic
variance of the IPW estimator relative to using the true propensity
score}, a result sometimes called the \emph{estimated propensity score
paradox} \citep{robins1986new}.

\end{example}

\begin{remark}[Sandwich variance estimator]
\label{w9:rem:sandwich}
Equation~\eqref{w9:eq:if-corrected} also provides the basis for the
\emph{sandwich variance estimator}.  Replacing $\theta_0$ by
$\hat\theta$ in all quantities, the asymptotic variance of $\hat\psi$
is consistently estimated by
\[
  \hat V
  =
  \frac{1}{n^2}\sum_{i=1}^n \hat\varphi_i^2,
\]
where $\hat\varphi_i$ is $\varphi(O_i)$ evaluated at
$(\hat\psi,\hat\alpha)$.  This estimator automatically accounts for
nuisance estimation uncertainty and is valid without any re-sampling.
Standard errors from a logistic regression routine that ignores the
link between $\hat\alpha$ and $\hat\psi$ will in general be
\emph{incorrect}.
\end{remark}

\begin{warn}{The central challenge in causal estimation}
The primary difficulty in practice is often not identifying the target
parameter, but \emph{correctly accounting for the effect of nuisance
estimation} on the final estimator.  Plugging a misspecified
$\hat\alpha$ into $U_1$ without the correction term in
\eqref{w9:eq:if-corrected} yields not only a potentially biased
estimate of $\psi$ but also an incorrect variance estimate.  The
Z-estimation framework resolves both issues simultaneously.
\end{warn}

\begin{remark}[Machine learning for nuisance estimation]
\label{w9:rem:ml}
The Z-estimation theorem as stated requires $\hat\alpha$ to converge
at rate $n^{-1/2}$, which holds for finite-dimensional parametric
models but fails for nonparametric or machine-learning estimators.
When flexible methods are used for $\pi(x)$ or $\mu_t(x)$, the
correction term in \eqref{w9:eq:if-corrected} no longer has the simple
form derived here, and a more careful analysis based on
\emph{sample splitting} or \emph{cross-fitting} is needed.  This is
the starting point for the doubly robust and debiased machine-learning
estimators introduced in Chapter~10.
\end{remark}

% ============================================================
\section{Variance Estimation and Confidence Intervals}
\label{w9:sec:variance}
% ============================================================

Once an estimator admits the asymptotic linear representation
\[
  \sqrt{n}(\hat\psi - \psi)
  =
  \frac{1}{\sqrt{n}}\sum_{i=1}^n \varphi(O_i) + o_p(1),
\]
its asymptotic variance is
\[
  \mathrm{Var}(\hat\psi)
  \approx
  \frac{1}{n}\,\mathrm{Var}\{\varphi(O)\}.
\]
A natural estimator of this variance is obtained by replacing
$\varphi(O_i)$ with estimated influence values $\hat\varphi_i$:
\[
  \hat V
  =
  \frac{1}{n^2}\sum_{i=1}^n \hat\varphi_i^2,
\]
or, equivalently,
\[
  \hat V
  =
  \frac{1}{n(n-1)}\sum_{i=1}^n (\hat\varphi_i - \bar\varphi)^2,
  \qquad
  \bar\varphi = \frac{1}{n}\sum_{i=1}^n \hat\varphi_i.
\]
This yields the approximate Wald confidence interval
\[
  \hat\psi \pm z_{1-\alpha/2}\,\sqrt{\hat V}.
\]

\begin{prop}{Variance from the influence function}
\label{w9:prop:variance}
If
\[
  \sqrt{n}(\hat\psi - \psi)
  =
  \frac{1}{\sqrt{n}}\sum_{i=1}^n \varphi(O_i) + o_p(1)
\]
with $\E\{\varphi(O)\} = 0$ and $\E\{\varphi(O)^2\} < \infty$, then
\[
  \sqrt{n}(\hat\psi - \psi)
  \overset{d}{\longrightarrow}
  N\!\left(0,\, \E[\varphi(O)^2]\right).
\]
Consequently,
\[
  \mathrm{Var}(\hat\psi) = \frac{1}{n}\,\E[\varphi(O)^2] + o(n^{-1}).
\]
\end{prop}

The proof follows immediately from the central limit theorem applied to
the i.i.d.\ sum.  The influence function thus plays a dual role: it
characterizes both the asymptotic distribution and the asymptotic
variance.

\begin{remark}[Estimating $\varphi(O_i)$ in practice]
\label{w9:rem:plugin-if}
In practice, $\varphi(O_i)$ depends on unknown parameters and nuisance
functions, so it is replaced by an estimated influence value
$\hat\varphi_i$.  This substitution is not innocuous: the plug-in
variance formula $\hat V = n^{-2}\sum_i\hat\varphi_i^2$ is valid only
when the estimator is asymptotically linear with influence function
$\varphi$ \emph{and} the error from replacing unknown quantities by
estimates is asymptotically negligible.  For the plug-in estimators
discussed in Section~\ref{w9:sec:nuisance}, these conditions require
additional regularity --- in particular, the nuisance estimators must
converge fast enough that the plug-in error is $o_p(n^{-1/2})$.  We
return to this point in Chapter~10, where doubly robust estimators
make the role of nuisance estimation in variance estimation explicit.
\end{remark}

% ============================================================
\section{Toward Efficiency: Semiparametric Models and the Efficient Influence Function}
\label{w9:sec:efficiency}
% ============================================================

Section~\ref{w9:sec:asymlin} showed that the influence function of an
asymptotically linear estimator determines both its asymptotic
distribution and its asymptotic variance.  Different regular estimators
of the same parameter may therefore have different large-sample
variances.  This raises a natural question: among all regular estimators
in a given statistical model, what is the smallest achievable asymptotic
variance?  Answering that question leads to semiparametric efficiency
theory.  In this section we introduce only the basic objects ---
semiparametric models, regular estimators, and the efficient influence
function --- as preparation for Chapter~10, where these ideas are used
to construct doubly robust and efficient estimators.  The geometric
interpretation of the efficiency bound and the systematic design of
estimators that achieve it are developed there.

\subsection{Semiparametric Models}
\label{w9:subsec:semiparam}

In causal inference, the parameter of interest $\psi = \Psi(P)$ is a
finite-dimensional quantity (e.g., the ATE), but the full distribution
$P$ of the observed data $O = (X, T, Y)$ is never fully specified.
This structure motivates the following definition.

\begin{defn}{Semiparametric Model}
\label{w9:defn:semiparam}
A \textbf{semiparametric model} is a statistical model $\mathcal{P}$
for the distribution $P$ of the observed data in which:
\begin{enumerate}[label=(\roman*)]
  \item the \textbf{parameter of interest} $\psi = \Psi(P) \in
    \mathbb{R}^p$ is finite-dimensional; and
  \item the remaining aspects of $P$ --- collectively called the
    \textbf{nuisance parameter} --- are infinite-dimensional (i.e.,
    $\mathcal{P}$ is not fully parametric).
\end{enumerate}
\end{defn}

The canonical example in this course is the nonparametric model
$\mathcal{P} = \{P : 0 < \pi(x) < 1\ \forall x\}$ with target
parameter $\psi = \E[Y(1)-Y(0)]$.  Here the nuisance consists of the
entire joint distribution of $(X, T, Y)$ subject to the positivity
constraint; no parametric form is assumed for $\mu_t(x)$ or $\pi(x)$.
In the ATE example, the nuisance parameter includes the outcome
regressions $\mu_t(x)$, the propensity score $\pi(x)$, and the
marginal distribution of $X$; these are precisely the objects that must
be estimated in practice --- and were treated as estimated nuisance
parameters in Section~\ref{w9:sec:nuisance} --- but are not themselves
the target of inference.

\begin{remark}[Semiparametric vs.\ parametric efficiency]
\label{w9:rem:vs-param}
In a \emph{fully parametric} model, the Cramér--Rao lower bound gives
the minimum variance of any unbiased estimator of $\psi$.  When the
nuisance is infinite-dimensional, the classical Cramér--Rao bound no
longer applies directly: there are infinitely many directions in which
the likelihood can vary, and the relevant lower bound must account for
all of them.  The semiparametric efficiency bound is the analogue of
the Cramér--Rao bound for this setting; see \citet{bickel1993} and
\citet{tsiatis2006semiparametric} for comprehensive treatments.
\end{remark}

\subsection{Regular Estimators}
\label{w9:subsec:regular}

The efficiency bound is meaningful only within a restricted class of
estimators.  Without such a restriction, pathological estimators can
behave well at a single distribution while behaving erratically under
nearby perturbations, making variance comparisons meaningless from a
local asymptotic point of view.  The relevant class is that of
\emph{regular} estimators, which, informally, are estimators whose
asymptotic distribution is stable under small perturbations of the
data-generating distribution.

\begin{defn}{Regular Estimator (informal)}
\label{w9:defn:regular}
An estimator $\hat\psi$ of $\psi_0 = \Psi(P_0)$ is called
\textbf{regular} at $P_0$ if, along every smooth parametric submodel
$\{P_t : t \in \mathbb{R}\}$ with $P_0 = P_{t=0}$, the limiting
distribution of $\sqrt{n}(\hat\psi - \psi_t)$ under $P_t = P_{1/\sqrt{n}}$
does not depend on the submodel or its direction.
\end{defn}

Regularity rules out pathological estimators that exploit specific
features of the data-generating mechanism in a non-uniform way
(so-called \emph{superefficient} estimators).  All of the estimators
encountered in this course --- regression estimators, IPW estimators,
and their augmented variants --- are regular.  The formal theory of
regular estimators and the convolution theorem underlying the
efficiency bound are developed in \citet{vandervaart1998}
(Chapters~8 and~25) and \citet{tsiatis2006semiparametric} (Chapter~4).

\subsection{The Semiparametric Efficiency Bound and the Efficient Influence Function}
\label{w9:subsec:eif}

The following theorem, due to \citet{bickel1993}, is the central result
of semiparametric efficiency theory.

\begin{thm}{Semiparametric Convolution Theorem}
\label{w9:thm:convolution}
In a semiparametric model $\mathcal{P}$, the asymptotic distribution of
any regular estimator $\hat\psi$ of $\psi_0$ is at least as dispersed
as $N(0, V^*)$ in the sense that
\[
  \sqrt{n}(\hat\psi - \psi_0) \overset{d}{\longrightarrow}
  N(0, V^*) * M
\]
for some distribution $M$, where $*$ denotes convolution and
\[
  V^* = \E_P\!\bigl[\varphi^*(O)\,\varphi^*(O)^\top\bigr]
\]
is the \textbf{semiparametric efficiency bound}.  The function
$\varphi^*(O)$ with $\E[\varphi^*(O)] = 0$ and $\E[\varphi^*(O)^2] =
V^*$ is called the \textbf{efficient influence function} (EIF).  An
estimator achieves $V^*$ if and only if it is asymptotically linear
with influence function $\varphi^*(O)$.
\end{thm}

The efficient influence function is not merely any valid influence
function; it is the unique one that achieves the smallest asymptotic
variance among regular estimators in the model.  Every regular
estimator has some influence function, but only the estimator whose
influence function equals $\varphi^*$ attains the bound $V^*$ with no
convolution remainder.

\begin{remark}[How the EIF is computed]
\label{w9:rem:eif-computation}
The EIF $\varphi^*(O)$ can be derived as the pathwise derivative of
$\Psi(P)$ in the \emph{least favorable} direction, i.e., the direction
in $\mathcal{P}$ that is hardest to distinguish from $P_0$.  In
practice, one constructs $\varphi^*(O)$ by: (i) computing the
pathwise derivative $\varphi(O; P_\epsilon)$ for an arbitrary
parametric submodel (as in Example~\ref{w9:ex:if-ols}), and (ii)
projecting onto the tangent space of $\mathcal{P}$ at $P_0$.  For
nonparametric or large semiparametric models the tangent space is rich
enough that this projection is often the identity, and the EIF equals
the ordinary pathwise derivative.  We will carry out this construction
explicitly for causal parameters in Chapter~10.
\end{remark}

\begin{defn}{Semiparametrically Efficient Estimator}
\label{w9:defn:efficiency}
A regular estimator $\hat\psi$ is called \textbf{semiparametrically
efficient} at $P_0 \in \mathcal{P}$ if
\[
  \sqrt{n}(\hat\psi - \psi_0)
  \overset{d}{\longrightarrow}
  N(0, V^*),
\]
i.e., it achieves the semiparametric efficiency bound $V^*$ with no
additional convolution component $M$.  Equivalently, $\hat\psi$ is
asymptotically linear with influence function $\varphi^*(O)$.
\end{defn}

\subsection{The Efficient Influence Function for the ATE}
\label{w9:subsec:eif-ate}

We record the EIF for the ATE as it will be derived and used
extensively in Chapter~10.

\begin{remark}[EIF for the ATE]
\label{w9:rem:eif-ate}
For the ATE $\psi = \E\{Y(1) - Y(0)\}$ in the nonparametric
observational-data model under strong ignorability and positivity, the
efficient influence function is
\[
  \varphi^*(O)
  =
  \frac{T\{Y - \mu_1(X)\}}{\pi(X)}
  - \frac{(1-T)\{Y - \mu_0(X)\}}{1-\pi(X)}
  + \mu_1(X) - \mu_0(X) - \psi.
\]
This expression has a natural decomposition: the first two terms are an
IPW-style residual correction, and the last two are the regression
estimator centered at $\psi$.  Notably, the EIF depends on \emph{both}
nuisance functions $\pi(X)$ and $\mu_t(X)$, and an estimator that
plugs in consistent estimates of both achieves the efficiency bound
$V^* = \E[\varphi^*(O)^2]$.
\end{remark}

\begin{remark}[Bridge to Chapter~10]
\label{w9:rem:bridge}
Two important properties of $\varphi^*(O)$ in Remark~\ref{w9:rem:eif-ate}
will be central in Chapter~10:
\begin{enumerate}[label=(\roman*)]
  \item \textbf{Double robustness.}  The estimator obtained by
    replacing $\pi(X)$ and $\mu_t(X)$ by estimates and averaging
    $\varphi^*(O_i)$ is consistent for $\psi$ if \emph{either}
    $\hat\pi$ or $\hat\mu_t$ is consistent — not necessarily both.
    This is the \emph{doubly robust} property.
  \item \textbf{Semiparametric efficiency.}  When both nuisance
    estimators are consistent and converge at suitable rates, the
    resulting estimator is asymptotically linear with influence function
    $\varphi^*(O)$ and therefore achieves the efficiency bound $V^*$.
\end{enumerate}
The \emph{augmented IPW} (AIPW) estimator, sometimes called the
one-step or debiased estimator, is the principal tool for exploiting
these two properties simultaneously, and will be derived and analyzed
in Chapter~10.  The efficient influence function for the ATE will
become central there, where we show that the same object generates
both doubly robust estimating equations and the semiparametric
efficiency bound.
\end{remark}

% ============================================================
\section{Summary}
\label{w9:sec:summary}
% ============================================================

This chapter introduced the statistical framework underlying estimation
and inference for causal parameters.

\begin{enumerate}
  \item \textbf{Identification is not estimation.}  Identification
    provides a population formula, but not automatically a satisfactory
    estimator.

  \item \textbf{Estimating equations.}  Many estimators are defined as
    solutions to sample moment conditions
    $\mathbb{P}_n\{U(O;\theta)\}=0$, with the population condition
    $\E\{U(O;\theta_0)\}=0$ ensuring consistency.

  \item \textbf{Asymptotic linearity.}  Estimating equations often
    yield a first-order expansion
    $\sqrt{n}(\hat\theta - \theta_0) = n^{-1/2}\sum_i\varphi(O_i) +
    o_p(1)$, which immediately implies asymptotic normality via the
    CLT.

  \item \textbf{Influence functions.}  The function $\varphi(O)$
    describes the first-order sensitivity of the estimator to a single
    observation.  It is the key object for both asymptotic theory and
    variance estimation.

  \item \textbf{Nuisance estimation.}  In causal inference, nuisance
    functions must be estimated, and this estimation error propagates
    into the target estimator.  Controlling this propagation is the
    central statistical challenge.

  \item \textbf{Variance from the influence function.}  The asymptotic
    variance is $n^{-1}\E[\varphi(O)^2]$, consistently estimated by
    $n^{-2}\sum_i\hat\varphi_i^2$.

  \item \textbf{Efficiency.}  Among regular estimators, the one with
    the smallest asymptotic variance is efficient.  The efficient
    influence function characterizes this lower bound and will guide
    estimator construction in Chapter~10.
\end{enumerate}

\begin{tcolorbox}[defstyle, title={Key Objects in This Chapter}]
\begin{center}
\renewcommand{\arraystretch}{1.5}
\begin{tabular}{p{4.5cm} p{8cm}}
\toprule
\textbf{Object} & \textbf{Role} \\
\midrule
Estimating function $U(O;\theta)$
  & Defines the estimator via
    $\mathbb{P}_n\{U(O;\hat\theta)\}=0$. \\
Influence function $\varphi(O)$
  & First-order term in the expansion of
    $\hat\psi - \psi$; used for CLT and variance estimation. \\
Asymptotic variance $\E[\varphi(O)^2]/n$
  & Governs precision; the target for efficiency comparisons. \\
Efficient influence function $\varphi^*(O)$
  & Achieves the semiparametric efficiency bound; basis for
    doubly robust estimators (Chapter~10). \\
\bottomrule
\end{tabular}
\end{center}
\renewcommand{\arraystretch}{1}
\end{tcolorbox}

% ============================================================
\section*{Problems}
\addcontentsline{toc}{section}{Problems}
% ============================================================

\begin{enumerate}

\item \textbf{Estimating equations and moment conditions.}
  \begin{enumerate}
    \item Let $O = (X, Y)$ and define $\theta = \mathrm{Cov}(X,Y) /
      \mathrm{Var}(X)$ (the coefficient in the population simple
      regression of $Y$ on $X$).  Write down an estimating equation
      for $\theta$ and verify that the population moment condition
      holds at $\theta_0$.
    \item Let $\theta = F^{-1}(0.5)$ be the population median.
      Propose an estimating function $U(O;\theta)$ and verify the
      population moment condition.  [\emph{Hint}: consider
      $U(O;\theta) = \mathbf{1}(Y \leq \theta) - 0.5$.]
    \item Show that the OLS estimator and the IPW estimator of the
      ATE (Example~\ref{w9:ex:ipw-ee}) are both special cases of the
      general M-estimator framework.
  \end{enumerate}

\item \textbf{Asymptotic linearity and the CLT.}
  \begin{enumerate}
    \item Let $\hat\psi = \bar Y$ be the sample mean of i.i.d.\
      $Y_1,\dots,Y_n$ with $\E Y = \psi$ and $\mathrm{Var}(Y) =
      \sigma^2 < \infty$.  Write down the influence function, state
      its asymptotic distribution, and give a consistent estimator of
      $\mathrm{Var}(\hat\psi)$.
    \item Suppose $\hat\psi_1$ and $\hat\psi_2$ are two asymptotically
      linear estimators of the same parameter with influence functions
      $\varphi_1$ and $\varphi_2$.  Show that
      $\hat\psi_\lambda = \lambda\hat\psi_1 + (1-\lambda)\hat\psi_2$
      is also asymptotically linear for any fixed $\lambda \in [0,1]$,
      and find its influence function.
    \item Use part~(b) to explain why, among linear combinations of
      two estimators, the one minimizing asymptotic variance is
      generally not the simple average ($\lambda = 1/2$).
  \end{enumerate}

\item \textbf{Influence functions for causal estimators.}
  Consider the ATE $\psi = \E\{Y(1) - Y(0)\}$ identified under
  strong ignorability.
  \begin{enumerate}
    \item Assume $\pi(X)$ and $\mu_t(X)$ are both known.  Compute the
      influence functions of:
      \begin{enumerate}
        \item[(i)] the regression estimator
          $\hat\psi_{\mathrm{reg}} = \mathbb{P}_n\{\mu_1(X) -
          \mu_0(X)\}$;
        \item[(ii)] the IPW estimator
          $\hat\psi_{\mathrm{IPW}} = \mathbb{P}_n\!\bigl\{TY/\pi(X) -
          (1-T)Y/\{1-\pi(X)\}\bigr\}$.
      \end{enumerate}
    \item Under what conditions on the data-generating process will
      $\hat\psi_{\mathrm{reg}}$ have a smaller asymptotic variance
      than $\hat\psi_{\mathrm{IPW}}$?
    \item The efficient influence function for the ATE is given in
      Remark~\ref{w9:rem:eif-ate}.  Verify that it has mean zero
      under $P$ by computing $\E[\varphi^*(O)]$.
  \end{enumerate}

\item \textbf{Variance estimation.}
  \begin{enumerate}
    \item Let $\hat\varphi_i = \hat\mu_1(X_i) - \hat\mu_0(X_i) -
      \hat\psi_{\mathrm{reg}}$ be the estimated influence values for
      the regression estimator.  Write down the plug-in variance
      estimator $\hat V$ and simplify.  Under what conditions is
      $\hat V$ consistent for $\mathrm{Var}(\hat\psi_{\mathrm{reg}})$?
    \item A researcher reports $\hat\psi = 2.4$ and $\hat V = 0.09$
      based on $n = 400$ observations.  Compute a 95\% Wald confidence
      interval.  Is the treatment effect statistically distinguishable
      from zero at the 5\% level?
    \item Explain why simply plugging in $\hat\mu_t$ and $\hat\pi$ into
      the influence function formula may yield an inconsistent variance
      estimator when these nuisance estimators converge at a slower
      rate than $n^{-1/2}$.
  \end{enumerate}

\item \textbf{Efficiency comparisons.}
  \begin{enumerate}
    \item Suppose $\varphi_{\mathrm{reg}}$ and
      $\varphi_{\mathrm{IPW}}$ denote the influence functions of the
      regression and IPW estimators, respectively.  Without
      calculation, explain why neither estimator is always more
      efficient than the other.
    \item The semiparametric efficiency bound for the ATE is
      $\E[\varphi^*(O)^2]$, where $\varphi^*$ is given in
      Remark~\ref{w9:rem:eif-ate}.  Show that
      $\E[\varphi^*(O)^2] \leq \E[\varphi_{\mathrm{IPW}}(O)^2]$ by
      expanding the squared EIF.  [\emph{Hint}: use the law of
      iterated expectations to show that the cross-terms cancel.]
    \item Give one practical reason why an efficient estimator based
      on the EIF may not always be preferred over a simpler, less
      efficient estimator.
  \end{enumerate}

\end{enumerate}

% Bibliography (standalone only)
\ifSubfilesClassLoaded{%
  \bibliographystyle{plainnat}
  \bibliography{causal_inference}
}{}

\end{document}